\documentclass[12pt]{article}
\usepackage{amsmath}
\usepackage{enumitem}
\usepackage{geometry}
\geometry{margin=1in}
\usepackage{hyperref}
\hypersetup{
    colorlinks=true,
    linkcolor=black,
    urlcolor=blue,
    pdfborder={0 0 0}
}
\title{Ujian Akhir Semester – Kriptografi}
\date{}
\begin{document}
\maketitle

\section*{Petunjuk}
\begin{itemize}
    \item Pilih jawaban terbaik untuk setiap soal.
    \item Lingkari huruf pilihan Anda (A, B, C, atau D).
    \item Semua soal bernilai sama.
\end{itemize}

\section*{Soal Pilihan Ganda}
\begin{enumerate}[label=\textbf{Soal \arabic*.}, leftmargin=*, itemsep=1.5em]

\item \textbf{True Random Number Generator (TRNG) vs. Pseudo‑Random Number Generator (PRNG)}. Manakah pernyataan yang \textbf{benar}?
\begin{enumerate}[label=\Alph*.]
    \item TRNG menghasilkan urutan deterministik dari seed.
    \item PRNG dapat menghasilkan bit yang tidak dapat diprediksi tanpa mengetahui seed.
    \item TRNG biasanya lebih lambat tetapi tidak dapat diprediksi.
    \item PRNG selalu aman untuk menghasilkan kunci kriptografi.
\end{enumerate}

\item Linear Congruential Generator (LCG) tidak cocok untuk kriptografi karena:
\begin{enumerate}[label=\Alph*.]
    \item Memiliki periode tak terbatas.
    \item Pola output dapat dipulihkan dari sedikit nilai output.
    \item Menggunakan operasi modular yang mahal.
    \item Membutuhkan sumber entropi tinggi.
\end{enumerate}

\item Pada \textit{Blum‑Blum‑Shub} (BBS) CSPRNG, keamanan utama bergantung pada:
\begin{enumerate}[label=\Alph*.]
    \item Kesulitan faktorisasi bilangan komposit $n = p q$.
    \item Kesulitan menghitung logaritma diskrit.
    \item Kualitas generator \texttt{rand()} standar.
    \item Kecepatan iterasi modular.
\end{enumerate}

\item Steganografi LSB pada gambar beresolusi $640\times480$ berwarna (24‑bit). Kapasitas maksimum pesan (dalam bit) yang dapat disembunyikan adalah:
\begin{enumerate}[label=\Alph*.]
    \item $640 \times 480 = 307\,200$ bit
    \item $640 \times 480 \times 3 = 921\,600$ bit
    \item $640 \times 480 \times 2 = 614\,400$ bit
    \item $640 \times 480 \times 4 = 1\,228\,800$ bit
\end{enumerate}

\item Teknik \textit{steganalysis} yang paling umum untuk mendeteksi penyisipan LSB pada gambar adalah:
\begin{enumerate}[label=\Alph*.]
    \item Analisis frekuensi warna RGB.
    \item Pengujian statistik distribusi bit LSB (chi‑square).
    \item Pencarian pola berulang pada header file.
    \item Penerapan transformasi Fourier.
\end{enumerate}

\item Algoritma kuantum yang mengancam keamanan RSA dan ECC adalah:
\begin{enumerate}[label=\Alph*.]
    \item Algoritma Grover.
    \item Algoritma Shor.
    \item Algoritma Simon.
    \item Algoritma Deutsch‑Jozsa.
\end{enumerate}

\item Post‑Quantum Cryptography (PQC) dibagi menjadi lima kategori utama. Manakah yang \textbf{bukan} termasuk dalam kategori tersebut?
\begin{enumerate}[label=\Alph*.]
    \item Lattice‑based cryptography
    \item Hash‑based signatures
    \item Code‑based cryptography
    \item Symmetric‑key cryptography
\end{enumerate}

\item Keamanan \textbf{lattice‑based cryptography} biasanya didasarkan pada kesulitan:
\begin{enumerate}[label=\Alph*.]
    \item Faktor prima besar.
    \item Shortest Vector Problem (SVP) pada lattice.
    \item Diskret logarithm pada kurva eliptik.
    \item Decoding kode linear umum.
\end{enumerate}

\item Pada problem \textit{Learning With Errors} (LWE) yang digunakan dalam protokol pertukaran kunci, apa yang diketahui oleh penyerang?
\begin{enumerate}[label=\Alph*.]
    \item Matriks publik $A$ dan vektor $b = A s + e$.
    \item Hanya matriks $A$.
    \item Hanya vektor $b$.
    \item Rahasia $s$ secara langsung.
\end{enumerate}

\item Kyber, salah satu finalist NIST PQC, menggunakan transformasi apa untuk perkalian polinomial?
\begin{enumerate}[label=\Alph*.]
    \item Fast Fourier Transform (FFT)
    \item Number Theoretic Transform (NTT)
    \item Karatsuba multiplication
    \item Montgomery reduction
\end{enumerate}

\item Pada \textbf{Lamport One‑Time Signature}, keamanan kunci privat dan publik dicapai dengan:
\begin{enumerate}[label=\Alph*.]
    \item Enkripsi RSA.
    \item Hashing nilai acak dengan SHA‑256.
    \item Eksponensiasi modular.
    \item Penggunaan fungsi PRF.
\end{enumerate}

\item XMSS (Extended Merkle Signature Scheme) mengandalkan struktur data apa untuk menghasilkan jalur otentikasi?
\begin{enumerate}[label=\Alph*.]
    \item Daftar berantai.
    \item Pohon Merkle.
    \item Graf berarah.
    \item Tabel hash.
\end{enumerate}

\item Keamanan \textbf{McEliece Cryptosystem} bergantung pada kesulitan:
\begin{enumerate}[label=\Alph*.]
    \item Faktorisasi bilangan komposit.
    \item Decoding kode linear umum (General Decoding Problem).
    \item Discrete logarithm pada kurva eliptik.
    \item Menemukan akar kubik modulo prima.
\end{enumerate}

\item Dalam skenario \textit{Hybrid TLS 1.3} yang menggabungkan ECDHE dengan Kyber, langkah pertama yang dilakukan oleh client adalah:
\begin{enumerate}[label=\Alph*.]
    \item Mengirim \texttt{ClientHello} berisi \texttt{KeyShare} ECDHE saja.
    \item Mengirim \texttt{ClientHello} berisi \texttt{KeyShare} Kyber saja.
    \item Mengirim \texttt{ClientHello} berisi \texttt{KeyShare} ECDHE \textbf{dan} paket Kyber.
    \item Menunggu \texttt{ServerHello} sebelum mengirim apapun.
\end{enumerate}

\item Menurut roadmap migrasi ke post‑quantum yang disarankan, jangka waktu mana yang paling tepat untuk memulai evaluasi algoritma PQC?
\begin{enumerate}[label=\Alph*.]
    \item 2024–2026
    \item 2026–2030
    \item 2030–2035
    \item 2035–2040
\end{enumerate}

\end{enumerate}

\newpage
\section*{Kunci Jawaban}
\begin{enumerate}[label=\textbf{Soal \arabic*.}]
    \item C
    \item B
    \item A
    \item B
    \item B
    \item B
    \item D
    \item D
    \item B
    \item A
    \item B
    \item B
    \item B
    \item B
    \item C
    \item A
\end{enumerate}

\end{document}
